\section{Introduction}
\label{sec:intro}

%TL;DR (X) Generating discrete 3D structures such as Minecraft buildings is useful and expensive, because volume grows cubically.
Block worlds such as Minecraft hold millions of player-made buildings, and generating new ones is a
standing goal for procedural content and for 3D generation research~\citep{salge2018gdmc,chen2019voxelcnn,merino2026dreamcubed}.
A building in such a world is a voxel grid in which every cell holds one of hundreds of block types.
Autoregressive models that write this grid cell by cell work~\citep{chen2019voxelcnn,merino2026dreamcubed},
but their sequences grow with the volume: a $32^3$ house spans up to 28{,}672 cells, and at $48^3$ the
largest houses in our corpus span 106{,}032.

%TL;DR (X) The exterior of a building is six images, and image-like 2.5D representations are the standard way to make 3D generation cheap.
Most of what a player sees of a building is its exterior: four walls, a roof and a footprint. Each of
these is an image, a face that records, for every cell, how deep the first block lies and which block
it is. Six such faces cost a quadratic, not cubic, number of tokens (at most 5{,}632 cells at $32^3$
and 13{,}442 at $48^3$ in our corpus). Continuous 3D generation already exploits the same idea:
multi-view and 2.5D methods generate depth or normal images and fuse them into a
shape~\citep{wang2024mvdd,lu2024direct25,long2024wonder3d,latents25d2025}.

%TL;DR (Y) The hard part is agreement: independently written faces describe different objects, and free decoding almost never produces a consistent set.
The difficulty is that six faces describe one object only if they agree. Where two faces see the same
voxel they must report the same block, and no face may report a first hit behind a voxel that another
face has already filled. A transformer that writes the six faces of a house in one sequence learns the
style of each face but not this agreement. Under free decoding, between 0.0\% and 1.8\% of its face
sets equal the projection of their own lift, and about three cells in ten change when we lift a face
set and project it back (\cref{tab:consistency}). The lift resolves every disagreement by overwriting,
and the result is a building in 34 pieces where a real house has three.

%TL;DR (Y) Continuous multi-view methods treat consistency as approximate and learn it; in a block world it is exact, so it can be enforced rather than learned.
Continuous multi-view generators face the same problem and address it with cross-view attention,
synchronized denoising or a learned fusion step~\citep{shi2024mvdream,liu2024syncdreamer,wang2024mvdd},
which reduces disagreement without removing it. A block world differs in one useful way: its faces are
exact projections of a discrete grid, so whether two faces agree is a yes-or-no question about
individual voxels. This makes consistency a constraint that a decoder can enforce, in the same way
that grammar-constrained decoding enforces syntax in language models~\citep{scholak2021picard,geng2023grammar,willard2023outlines}.

%TL;DR (Z) Our method: cross-face constrained decoding keeps a voxel state of what earlier faces committed and masks the logits of later faces to what agrees.
We introduce \cfd{} decoding. While the model writes the faces in order, a shell state records every
voxel that an earlier face has committed: the first block along each ray as known, and the cells in
front of it as empty. For each cell of a later face, the state determines which depths are still
possible and which block, if any, a depth must carry; the decoder masks the model's logits to exactly
that set. The model keeps every choice the earlier faces left open. The method needs no training, no
change to the model, and about 12\% extra sampling time.

%TL;DR (1a) Results: consistency becomes exact, fragmentation and render distance halve, it stacks with best-of-N, and the face generator matches a voxel transformer.
On 399 held-out Minecraft houses, \cfd{} decoding makes every face set consistent and writes 21\% to
25\% of all cells without sampling, because only one answer agrees with the faces before them. It
reduces the number of disconnected parts from 33.7 to 12.9, the multi-view render distance
\mvkid{} from 0.299 to 0.166, and \blockscore{}~\citep{anonymous2027blindness} from 679 to 589,
against a control decoder that enforces only per-face syntax. Combined with best-of-4 selection it
reaches an \mvkid{} of 0.098, the best of any learned generator we evaluate, at 7.7 parts per
building. In a paired comparison the face generator ties a 3D voxel transformer on geometry and beats it
on appearance (\cref{sec:vsvoxel}).

%TL;DR (1b) A dimensional study from one face to six faces to voxels shows the failure sits at the agreement between views, not at single-facade realism.
To locate where generating through 2D images breaks, we train the same model family on one facade
without depth, one facade with depth, and six faces, and compare each with a voxel model projected to
the same representation. \UNMEASURED{Single facades are close to the real distribution at every model size,
and the largest loss appears when six faces must agree, which is the loss \cfd{} decoding removes
(\cref{sec:ladder}).}

%TL;DR (2) Contributions as bullets.
In summary, we make the following contributions.
\begin{itemize}[leftmargin=*,itemsep=1pt,topsep=2pt]
  \item A face representation of block-building exteriors, six (depth, block) images with an exact
        lift and projection, and a measure of how far a face set is from consistent (\cref{sec:background}).
  \item \cfd{} decoding, a training-free sampler under which every generated face set is consistent
        by construction (\cref{sec:method}, \cref{prop:consistent}).
  \item Evidence that consistency, not model capacity, limited face generators: the constraint
        halves fragmentation and render distance, composes with best-of-$N$ selection, and lets a face
        generator match a voxel transformer on geometry at a quadratic token cost (\cref{sec:results}).
  \item A dimensional study from one facade to six faces to voxels that shows where 2D generation of
        3D structures breaks (\cref{sec:ladder}).
\end{itemize}

%TL;DR Future work in one line: larger grids, entrance-aware faces, other exact 2.5D domains such as LEGO.
Because faces scale with area rather than volume, the method points to larger buildings than voxel
models reach today; we discuss larger grids, entrance-aware ordering and other exact 2.5D domains in
\cref{sec:conclusion}.
